% ============================================================================
% Fichier  : partie3/P03_C09_cadre_juridique.tex
% Titre    : Le Cadre Juridique
% Auteur   : MINKA MI NGUIDJOÏ Thierry Emmanuel
% ============================================================================
% Sources : loi 2010/012, loi 2010/013, loi 2024/017 (Cameroun) ; Convention
% de Malabo (UA, 2014) ; Convention de Budapest (2001).
% ============================================================================

\chapter{Le Cadre Juridique}
\label{chap:droit}
\chaptermark{Le cadre juridique}

\epigraph{%
    « Les journaux tracent la frontière entre la faute professionnelle\\
    et la bonne foi prouvée. »%
}{Principe du cours}

\niveauChapitre{Fondamental}{Chapitre~\ref{chap:normes}~; notions de droit
numérique.}

\begin{boxobjectifs}
\begin{itemize}
    \item Connaître la loi camerounaise 2010/012 et la loi 2010/013.
    \item Situer la loi de 2024 sur les données personnelles.
    \item Comprendre le rôle et les pouvoirs de l'ANTIC.
    \item Articuler le droit national avec la Convention de Malabo et celle
          de Budapest.
    \item Déduire les obligations de LiZbiZ en cas d'incident.
\end{itemize}
\end{boxobjectifs}

\bigskip

% ============================================================================
\section*{Question directrice : qu'est-ce que la loi impose à LiZbiZ le 6 mars au matin~?}
% ============================================================================

L'incident n'est pas qu'un problème technique~: il déclenche des obligations.
Conserver les traces, signaler à l'ANTIC, notifier une fuite de données~:
autant de devoirs dont le non-respect engage la responsabilité. Ce chapitre
pose le cadre.

\medskip
\textbf{Plan du chapitre.}
\begin{enumerate}
    \item La loi 2010/012 et la loi 2010/013 (section~\ref{sec:c09-lois}).
    \item Les données personnelles (section~\ref{sec:c09-donnees}).
    \item L'ANTIC (section~\ref{sec:c09-antic}).
    \item Malabo et Budapest (section~\ref{sec:c09-international}).
    \item Les obligations de LiZbiZ (section~\ref{sec:c09-obligations}).
\end{enumerate}

\separateur

% ============================================================================
\section{La Loi 2010/012 et la Loi 2010/013}
\label{sec:c09-lois}
% ============================================================================

\begin{boxjuris}
La \loicam{loi n°~2010/012 du 21~décembre~2010} relative à la cybersécurité et
à la cybercriminalité est le texte fondateur~: elle réprime l'accès et le
maintien frauduleux dans un système, impose la conservation des traces et la
coopération avec l'ANTIC. La \loicam{loi n°~2010/013} régit les
communications électroniques et les obligations des opérateurs. Les articles
précis se lisent dans le texte publié au Journal officiel~; le cours n'en
cite que des éléments vérifiés.
\end{boxjuris}

% ============================================================================
\section{Les Données Personnelles}
\label{sec:c09-donnees}
% ============================================================================

\begin{boxjuris}
La \loicam{loi n°~2024/017 du 23~décembre~2024} régit la protection des
données à caractère personnel. Une fuite de données clients, comme celle
redoutée dans l'opération GAROUA, relève de ce texte~: information de
l'autorité, des personnes concernées, et documentation de la violation.
\end{boxjuris}

% ============================================================================
\section{L'ANTIC}
\label{sec:c09-antic}
% ============================================================================

\begin{boxdef}
\termecle{ANTIC} (Agence Nationale des Technologies de l'Information et de la
Communication)~: autorité camerounaise de la sécurité électronique. Elle
reçoit les signalements d'incidents, peut auditer, coordonne la réponse
nationale et coopère avec le parquet.
\end{boxdef}

% ============================================================================
\section{Malabo et Budapest}
\label{sec:c09-international}
% ============================================================================

\begin{tableaucours}{L{3.6cm}L{8.6cm}}{%
    \tableauHeader{Instrument} &
    \tableauHeader{Portée}
}
Convention de Malabo (UA, 2014) & Cadre continental~: cybersécurité, cybercriminalité, données personnelles \\
Convention de Budapest (2001) & Cadre international~: infractions, preuve électronique, coopération \\
\end{tableaucours}
\vspace{-0.8em}
\begin{center}{\small\itshape Tableau~9.1. Deux instruments supranationaux~\cite{malabo2014,budapest2001}}\end{center}

% ============================================================================
\section{Les Obligations de LiZbiZ}
\label{sec:c09-obligations}
% ============================================================================

\begin{boxscenario}{Opération GAROUA. Les devoirs du 6 mars}
\textbf{Constatation.} LiZbiZ subit un rançongiciel avec exfiltration
probable de données clients. Elle doit~: conserver les journaux, signaler
l'incident à l'ANTIC, documenter la violation, et, si des données
personnelles sont touchées, en informer l'autorité et les personnes.

\medskip
\textbf{Analyse} (confiance élevée). Ne pas disposer de journaux exploitables
affaiblirait la position de LiZbiZ~: la preuve de sa diligence repose sur les
traces. Le droit transforme la supervision en obligation, et le forensique en
moyen de preuve.
\end{boxscenario}

\begin{boxmetier}{RSSI et juriste TIC}
\textbf{Connaître ses obligations avant l'incident.} Le RSSI prépare, à
froid, la liste des signalements et notifications dus, avec leurs
destinataires et délais. En pleine crise, on applique une procédure, on ne
découvre pas le droit.
\end{boxmetier}

\separateur

\begin{boxresume}
\textbf{Ce qu'il faut retenir de ce chapitre~:}
\begin{itemize}
    \item \textbf{Loi 2010/012}~: accès frauduleux, conservation, coopération.
    \item \textbf{Loi 2024/017}~: protection des données personnelles.
    \item \textbf{L'ANTIC} reçoit les signalements et coordonne.
    \item \textbf{Malabo et Budapest} encadrent au niveau continental et
          international.
    \item \textbf{Les journaux} fondent la preuve de la diligence.
\end{itemize}
\end{boxresume}

\begin{boxinfo}
\textbf{Points de contrôle de l'élève.} Je sais~: citer les lois 2010/012,
2010/013, 2024/017~; décrire l'ANTIC~; situer Malabo et Budapest~; énoncer
les obligations de LiZbiZ.
\end{boxinfo}

\bigskip

\begin{boxexo}[title={\ding{45}\quad Exercice~9.1. La loi fondatrice \niveaubase}]
Quel texte camerounais fonde le droit de la cybersécurité~? De quelle année~?
\end{boxexo}

\begin{boxexo}[title={\ding{45}\quad Exercice~9.2. L'ANTIC \niveaubase}]
Citez trois pouvoirs de l'ANTIC.
\end{boxexo}

\begin{boxexo}[title={\ding{45}\quad Exercice~9.3. Données \niveaubase}]
Quel texte protège les données personnelles au Cameroun~?
\end{boxexo}

\begin{boxexo}[title={\ding{45}\quad Exercice~9.4. Malabo \niveaumoyen}]
Qu'apporte la Convention de Malabo par rapport au droit national~?
\end{boxexo}

\begin{boxexo}[title={\ding{45}\quad Exercice~9.5. Preuve de diligence \niveaumoyen}]
Pourquoi les journaux tracent-ils la frontière entre faute et bonne foi~?
\end{boxexo}

\begin{boxexo}[title={\ding{45}\quad Exercice~9.6. Notifications \niveaumoyen}]
Quelles notifications LiZbiZ doit-elle envisager le 6~mars~?
\end{boxexo}

\begin{boxexo}[title={\ding{45}\quad Exercice~9.7. Budapest \niveaumoyen}]
En quoi la Convention de Budapest facilite-t-elle la coopération
internationale~?
\end{boxexo}

\begin{boxexo}[title={\ding{45}\quad Exercice~9.8. À froid \niveaumoyen}]
Pourquoi préparer la liste des obligations avant l'incident~?
\end{boxexo}

\begin{boxexo}[title={\ding{45}\quad Exercice~9.9. Procédure de signalement \niveauexpert}]
Rédigez la procédure de signalement à l'ANTIC pour LiZbiZ~: déclencheur,
contenu, délai, destinataire.
\end{boxexo}

\begin{boxexo}[title={\ding{45}\quad Exercice~9.10. Note juridique \niveauexpert}]
Rédigez une note qui récapitule les obligations légales de LiZbiZ en cas
d'incident majeur.
\end{boxexo}

\bigskip

\begin{boxinfo}
\textbf{Pour aller plus loin.}
\begin{itemize}
    \item Lois 2010/012~\cite{loi2010012}, 2010/013~\cite{loi2010013},
          2024/017~\cite{loi2024017}.
    \item Convention de Malabo~\cite{malabo2014}~; Convention de
          Budapest~\cite{budapest2001}.
\end{itemize}
\end{boxinfo}

\begin{boxinfo}
\textbf{Transition.} Normes et droit convergent sur une question~: une trace
est-elle recevable~? C'est l'objet du Chapitre~\ref{chap:preuve}.
\end{boxinfo}
